\pagenumbering{roman}
\begin{titlepage}
\raggedright
{\sffamily\bfseries\Huge Methods Digest\par}
\vspace{4mm}
{\sffamily\LARGE How harms of privacy and data practices in health apps are communicated: the methodology of the reading corpus, stage by stage\par}
\vspace{10mm}
\ifdefined\onlystage
  {\sffamily\large Stage volume \textcolor{Stage\onlystage}{\textbf{\onlystage}}\par}
  \vspace{2mm}
  {\small This volume contains one stage chapter, the cards of the papers that belong to it, their references and their full texts. The complete digest is a separate file.\par}
\else
  {\sffamily\large Complete volume: nine stage chapters, cards for 99 papers, numbered references and 99 appended full texts\par}
\fi
\vspace{10mm}
\noindent\rule{\textwidth}{0.4pt}\par\medskip
{\small Compiled for the bachelor thesis of Hamza from the reading corpus in \texttt{01\_Literature}. This document is an organisational aid: it adds no thesis content, reports no results of the thesis and resolves no open decision. Methodological choices are recorded in \texttt{00\_Admin/Decision-Log.md}; this digest only collects the evidence for them.\par}
\vfill
{\small\color{black!70}\sffamily
\textbf{Private use.} The appended full texts include copies that may not be redistributed (for example Krippendorff 2019, a private copy). Do not publish or share this PDF outside the thesis project.\par}
\end{titlepage}

\tableofcontents
\clearpage
\pagenumbering{arabic}

\chapter*{How to use this digest}
\addcontentsline{toc}{chapter}{How to use this digest}
\markboth{How to use this digest}{How to use this digest}

\section*{The idea}
The thesis reading list contains about one hundred papers that matter for the methodology: how app reviews select and capture apps, how walkthroughs and codebooks are built, how two coders are compared, how language models are validated, how harms are modelled and how all of this is reported. This digest replaces re-reading them. Each stage of the planned study has a chapter that compares the approaches found in the literature, with the numbers, the design choices and the cautions; each paper has a card; and every citation is a link to the exact sentence in the paper.

\section*{Links}
\begin{itemize}
\item A reference like \textcolor{LinkBlue}{[12, p.\,4]} is a link. It jumps to page 4 of the appended copy of paper 12, where the cited sentence is highlighted. Page numbers are those of the appended copy (see the page-number remarks in the cards), not necessarily the journal pagination. In a PDF viewer use ``back'' (Alt+Left) to return.
\item The highlighted sentence carries a small orange mark; clicking it jumps back to the sentence that cites it.
\item An author--year label such as \emph{Malki et al.\ (2024) [1]} links to the numbered reference and to the paper's first page.
\item The bookmarks (outline panel) list the parts, the stage chapters, all cards and all full texts with their section headings.
\item In a stage volume, papers that belong to other stages are not appended; their labels stay as plain text and their numbers are the same as in the complete volume.
\end{itemize}

\section*{The nine stages}
\begin{table}[h]
\centering\small
\begin{tabularx}{\textwidth}{@{}c P{3.4cm} Y P{2.2cm}@{}}
\toprule
 & \textbf{Stage} & \textbf{Question answered} & \textbf{Decisions}\\
\midrule
\stagetag{S} & Scope and definitions & What is a tracker app, what is harm, privacy and transparency, what kind of study is this? & D3, D6, D11\\
\stagetag{A} & App search and selection & Which stores, keywords, ranking and thresholds; how is the sample documented? & D1--D5, D7\\
\stagetag{D} & Material capture & What is collected for each app and how is it archived and made citable? & D2, D5, D7\\
\stagetag{W} & Walkthrough & How are tasks, questions and outcomes defined for the interface inspection? & D10, D11\\
\stagetag{C} & Codebook and pilot & Where do the codes come from, what does an entry contain, how is the pilot run? & D8, D10\\
\stagetag{I} & Two-coder comparison & Which agreement statistic, which thresholds, which sample, how reported? & D8\\
\stagetag{L} & LLM feasibility & How is a language model configured, validated and reported? & D9\\
\stagetag{H} & Harm tree & How are harms defined, typed and organised into trees? & D12\\
\stagetag{R} & Reporting & Which reporting guidelines and checklists apply? & ---\\
\bottomrule
\end{tabularx}
\end{table}

\section*{Status and caveats}
\begin{itemize}
\item 100 papers are registered; 99 have a full text appended (Parker et al., no.\,41, has none) and all 99 have a card.
\item Nissenbaum 2009 (no.\,70) is only the table of contents and introduction, from a scanned copy with an OCR text layer; quotations follow the OCR text and may contain recognition errors.
\item Some appended files are author manuscripts, web captures or supplementary appendices rather than the typeset article; each card says so.
\item Card critiques are marked ``card''. They are the compiler's reading and should be checked against the paper before being used as the thesis' own position.
\item The thesis draft in \texttt{03\_Thesis} and the supervisor meeting of 13 October supersede the proposal where they conflict; the Decision Log lists twelve decisions that are still open.
\end{itemize}
